\section{An even affine horizon: a quartic in 209 unrestricted
integers}\label{an-even-affine-horizon-a-quartic-in-209-unrestricted-integers}

\subsection{Result, comparison and exact
source}\label{result-comparison-and-exact-source}

The equation in \nolinkurl{integer209/equation.json} is the sum of the
squares of 149 literal quadratic residuals in 209 independently
unrestricted integer variables. Its collected polynomial has degree 4,
5,486 nonzero terms and coefficient height 17592186044416=2\^{}44. The
maximum residual coefficient is 4194304=2\^{}22. Every coefficient is
supplied, with exact maps and an independent checker.

The immediate parent is the newly inspected Pell-growth210 equation from
snapshot \nolinkurl{c8187221a60df5baa0047c254a592a014491b60c}. Its original archive
is retained unchanged as \nolinkurl{inherited/Source210.zip}, SHA256
\nolinkurl{33fecaaa728d5f38608b8b0930c09b08dbe9a01e03f6269f092ec6142d72a0ee}. The
exact parent stage is \nolinkurl{equations/principal210/equation.json},
SHA256 \nolinkurl{20a2108353a3756ed7d1ddfdc71695131dcf0b6d94039988d9565ac48e6e1f0a}.
It is copied byte-for-byte as \nolinkurl{sources/parent210.json}.

The parent has 210 variables, 150 rows, 5,282 terms and height 2\^{}40.
This candidate saves one variable and one row, at the cost of 204
additional terms and sixteen times the height. It is a
dimension/coefficient tradeoff, not dominance. The separately verified
212-variable scaled-pair construction is unchanged and retained as a
method control. No global minimum or historical-priority claim is made,
and this bounded candidate is the final construction tested here.

The new ordinary proof, composed with the identified inherited
calibration, establishes in ZFC that this polynomial has an integer zero
exactly when there is a finite syntactic ZFC refutation. Negating gives
nonsolvability iff Con(ZFC). It does not prove consistency,
inconsistency or a zero. The argument is internal to ZFC, without a
restriction to externally standard proof codes.

The credited incoming source supplies standalone Pell growth, its
sharper constructor majorant, the affine radix, dyadic residue
permutation, squarefree beta spacing, exact projections and all logical
calibration. Those are not rediscoveries. The new contribution is the
even-affine horizon together with the signed bootstrap,
smaller-threshold estimates and changed cofinal parity construction
needed to remove its private root. Every unchanged row is literally
retained after one affine substitution and renumbering. No fresh
Lean/Rocq build is claimed; the full original ordinary-proof package is
retained, not represented as a complete compilable Lean dependency
closure.

\subsection{The literal change and its semantic
scope}\label{the-literal-change-and-its-semantic-scope}

Let delta=2\^{}19=524288. The incoming construction has the actual
horizon

\begin{verbatim}
t=16+4h_old^2.
\end{verbatim}

Replace it by

\begin{verbatim}
t=4h+2,             h independently unrestricted.          (1)
\end{verbatim}

The old root is coordinate 178 of the frozen 210-coordinate source and
occurs only in residual 125. Substitute 4h+2 for coordinate 0 in every
other residual, delete row 125 and coordinate 178, and renumber the
retained coordinates. Every residual stays quadratic. Exactly five
retained rows change: the relevant front endpoint row, shared frame row,
orientation row, standalone Pell product row and growth-divisibility
row. The certificate gives their full old/new coefficient arrays, not
just these descriptions.

The inherited reduction also commutes with this relaxation: throughout
all six projection stages the old root stays private; the five radical
directions and restoration polynomials avoid both t and the root; the
final affine consequence uses only the local front rows 5 and 18. All 79
final invertible operations avoid the horizon row/root and never change
t. Thus applying the semantic change to this authenticated optimized
parent does not conceal a previous use of the old horizon in another
retained condition.

This is a semantic zero-existence change, not a bijection between
arbitrary old witness tuples: the old horizons are 0 modulo 4, the new
horizons 2 modulo 4. The retained physical-coordinate map sends new h to
actual t=4h+2 and leaves all other retained physical coordinates
unchanged. On physical tuples with t=2 modulo 4, its inverse is
h=(t-2)/4. The omitted old root is not reconstructed. The two directions
of the consistency equivalence are proved below by arithmetic extraction
and a fresh cofinal witness construction.

Keep the original standalone center 4t+1, with existing coordinates F
and three private growth coordinates X,Z,u:

\begin{verbatim}
Z=tF,
X^2-16Z^2-8ZF=1,
F=delta*t*(2u+1).                                      (2)
\end{verbatim}

No enlargement to center 8t+1 is required. Keep the even representative
C, selector, affine radix and exact-power request

\begin{verbatim}
B*eta=C+F+1,   D_radix=2(C+F+1),   q=D_radix^x.          (3)
\end{verbatim}

The last equality is proof notation for its retained quadratic
certificate. All fifteen remaining requests, forty-five
base/index/output ports, thirteen principal pairs, anchor comparison
rows and normalized digit rows survive. The old square frame is retained
exactly, including a harmless duplicate t atom left over from the
obsolete growth request. Its34 actual affine atoms are fully
reconstructed in \nolinkurl{frame_certificate.json} before and after
substitution. Thus its equation remains W=2+sum(atom\^{}2), and every
actual requested base, index and output has square less than W.

\subsection{Soundness: recover the horizon before using positive
growth}\label{soundness-recover-the-horizon-before-using-positive-growth}

From (1), t is never 0 or -1 and \textbar t\textbar\textgreater=2. The
last row of (2) alone gives

\begin{verbatim}
F!=0, F even, |F|>=delta*|t|>=2^20.                    (4)
\end{verbatim}

No Pell classification at a positive growth center is used here. In
particular the actual frame contains F and has
W\textgreater=2+F\^{}2\textgreater258 even when t or F is negative. The
entire inherited signed C-anchor and orientation proof therefore applies
before any sign recovery. Briefly, its nonnegative cap forces a nonzero
anchor ordinate with absolute index\textbar C\textbar; squared Pell
ordinate divisibility and the comparison congruences identify that index
without assuming C positive. The other classified absolute indices are
strictly smaller than\textbar C\textbar. Reduction of the pooled
orientation modulo A-1 then gives a nonnegative residue smaller than
\textbar A-1\textbar, forcing every principal pair to be its canonical
signed pair. This uses only the retained actual frame bounds and cap,
which are authenticated coefficient by coefficient. Consequently the
inherited signed-power lemmas are available: a nonunit base forces a
nonnegative index and the genuine integer power.

Retain y=F(2v+1), so \textbar y\textbar\textgreater=\textbar F\textbar.
Because C and F are even, C+F+1 is odd; D\_radix is nonzero even and has
absolute value at least 2. Its power request first gives
x\textgreater=0. If x=0, the last endpoint

\begin{verbatim}
x=1+[1+(t+1)y]ut
\end{verbatim}

would make 1+(t+1)y divide -1. But t+1!=0 and

\begin{verbatim}
|1+(t+1)y|>=|t+1||y|-1>=|y|-1>1.
\end{verbatim}

Thus x\textgreater=1 and \textbar q\textbar\textgreater=2. The q-pool is
now genuinely nonunit. Its retained t and y+1 requests imply
t\textgreater=0 and y+1\textgreater=0. Since t=2 modulo 4,
t\textgreater=2. A negative y would be at most
-\textbar F\textbar\textless-1, so

\begin{verbatim}
t>=2,             y>=|F|>0.                          (5)
\end{verbatim}

This order is essential: positivity of t was not assumed to exclude x=0.
The rejected odd-affine horizon allowed t=-1,x=0,q=1 and failed exactly
here; \nolinkurl{review/REJECTED_ODD_HORIZON.md} preserves that
diagnostic.

Only now classify the standalone Pell equation at its positive odd
center 4t+1. The growth estimate below gives
\textbar F\textbar\textgreater t. If x=1, the first endpoint
x=F+1+(y+1)u0 would make the nonzero F divisible by y+1 while
\textbar F\textbar\textless=y\textless y+1, impossible. Thus
x\textgreater=2. The last endpoint now has a positive multiplier and
ut\textgreater=1, giving x\textgreater=2+(t+1)y. Hence
\textbar q+1\textbar\textgreater=3, so its pool is legitimate. Its F+t
request gives F+t\textgreater=0; together
with\textbar F\textbar\textgreater t this forces F\textgreater0. The C
request gives C\textgreater=0 and the already established anchor
nondegeneracy gives C\textgreater0. Thus
D\_radix\textgreater0,q\textgreater0, and the first endpoint gives
u0\textgreater=0 and x\textgreater=F+1. No postponed arithmetic sign is
used as its own premise.

\subsection{The original center and constructor bound work at
t=2}\label{the-original-center-and-constructor-bound-work-at-t2}

Write X\_n(a),Y\_n(a) for the standard Pell sequences, with initial
pair(1,0) and recurrence X\_next=aX+(a\^{}2-1)Y, Y\_next=X+aY. For
a\textgreater=2, every nonzero integer Pell ordinate has absolute value
Y\_n(a) for a positive integer n. The identities

\begin{verbatim}
Y_n(a)=n mod(a-1),
Y_(2n)(a)=2X_n(a)Y_n(a)
\end{verbatim}

give v2(Y\_n(a))=v2(n) when a is odd: X\_n is always odd and Y\_n is odd
exactly when n is odd. Also Y\_n(a)\textgreater=(2a-1)\^{}(n-1).

For the actual square-shell constructors, the largest value on arguments
at most M satisfies

\begin{verbatim}
E(M,M)+1=8M^2+14M+8<=8(M+1)^2.
\end{verbatim}

At t\textgreater=2, every local parameter and every constant leaf 0 or 1
has value-plus-one at most t. The unchanged authenticated maximum
constructor depth is 19. Induction therefore bounds every formula label
strictly below

\begin{verbatim}
L=8^(delta-1)*t^delta.                               (6)
\end{verbatim}

This sharper bound needs only t\textgreater=2, not t\textgreater=16. It
covers the same exact fourteen outputs and nine schema trees as the
preserved source ledger.

Equation(2) is X\^{}2-((4t+1)\^{}2-1)F\^{}2=1. Let
\textbar F\textbar=Y\_n(4t+1). Because 4t divides F, the first Pell
congruence gives 4t\textbar n.~The odd multiplier in(2) and the exact
valuation give v2(n)=19+v2(t). Hence n/(4t) is a positive integer of
exact valuation 17, and n\textgreater=delta*t. Consequently

\begin{verbatim}
|F|>=(8t+1)^(delta*t-1)>=(8t)^(delta*t-1)>=L^t.        (7)
\end{verbatim}

The exact ratio of the last lower bound to L\^{}t is
8\^{}(t-1)/t\textgreater=1. In particular
\textbar F\textbar\textgreater t, as used only after (5) in the
bootstrap. After F positivity has been recovered, (7) is the
weighted-label separation budget F\textgreater=L\^{}t.

Two other small-threshold arguments must be updated explicitly. Since
t\textgreater=2, 8t+1\textgreater=17 and delta*t-1\textgreater=t,
equation (7)'s first bound gives

\begin{verbatim}
F>=17^t>16^(t-1)>=lcm(1,...,t)^2.                    (8)
\end{verbatim}

Thus the all-primes binomial selector remains valid at t=2. We do not
reuse the false small-threshold assertion t\^{}t\textgreater=16\^{}t.
The same strong Pell bound also gives F\textgreater2\^{}32:
delta*t-1\textgreater=8 and 17\^{}8\textgreater2\^{}32. From
y\textgreater=F\textgreater t, t\textgreater=2 and
x\textgreater=2+(t+1)y, R0=2+t+x+y+F\textless2x. If
2\^{}s\textless=x\textless2\^{}(s+1), then s\textgreater=32 and

\begin{verbatim}
2^23 log2(2x)<2^23(s+2)<2^s<=x.
\end{verbatim}

The last inequality starts at s=32 and persists by induction. The
inherited literal local-degree bound therefore gives
R0\^{}(2\^{}23)\textless2\^{}x\textless=p\^{}(ex), so the same selected
local congruences recover exact local rows.

\subsection{Complete arithmetic extraction and logical
calibration}\label{complete-arithmetic-extraction-and-logical-calibration}

The unchanged C-cap makes every other nonnegative nominal index strictly
smaller than C. In particular F+t\textless C, all ordinary upper
arguments are at most C, and every needed lower digit index is at most
x. Equation (3) gives

\begin{verbatim}
D_radix=2(C+F+1)>2C,     q=D_radix^x>(2C)^x.
\end{verbatim}

Every binomial coefficient in a retained digit is therefore at most
C\^{}x, which is less than q/2\^{}x. The geometric-series tail is less
than its digit denominator divided by 32 because
x\textgreater=F+1\textgreater2\^{}32. The common normalized remainder
norm bounds the competing remainder in the centered half interval. Their
congruence thus forces equality. For the first digit, the nonzero
selector quotient gives
\textbar B\textbar\textless=C+F+1=D\_radix/2\textless q/2. Its true
positive binomial coefficient is less than q/32. The digit congruence
yields

\begin{verbatim}
B=binom(F+t,t)>0.
\end{verbatim}

The unchanged B-base power gives N=B\^{}x, while B\textbar D\_radix from
(3) gives N\textbar q. Thus all ten digit/range predicates, every local
interval, and all local precision assumptions are authenticated exactly
as in the incoming source. The all-primes selector (8) covers each row
and every prime of B in both soundness and completeness.

For labels bounded by L, equation (7) supplies the source's sharp
product separation theorem, including equality F=L\^{}t. Unequal factor
count s are impossible because F+1\textgreater L\^{}t. At equal count
m\textless=t, each positive elementary symmetric coefficient lies
in{[}1,L\^{}m{]} subset{[}1,F{]}, so differences have absolute value
less than F. The least nonzero coefficient cannot disappear on
evaluation at F. Thus label multisets coincide. Zero labels contribute
units and are omitted. The unchanged boundary and positive-accumulator
arguments recover a finite single-conclusion Horn flow; the inherited
closure proof extracts a refutation even if the flow has cycles. In
particular any exclusion of y=F is made after local recovery, not
before.

The literal local arithmetic uses the square-shell presentation. First
apply the retained injective square-shell/original-C222 recoding, which
carries finite refutations in both directions. This unchanged recoding
precedes the logical calibration; the shell labels are not silently
identified with the original numerical codes. The retained named
endpoints are \nolinkurl{KEndpoint.c222_iff_zf_refutation} and
\nolinkurl{ConstructibleChoice.zf_refutation_iff_zfc_refutation}. They
compose the literal arithmetic extraction with a finite ZFC refutation.
No machine execution or Jones-polynomial specialization is substituted
for these calibrated endpoints. This proves soundness.

\subsection{Cofinal completeness with the changed dyadic
bit}\label{cofinal-completeness-with-the-changed-dyadic-bit}

Start with a finite C222 refutation and all its finite local parameters,
proof length and padding. Choose a\textgreater=11 sufficiently large and
set

\begin{verbatim}
t=2^(2a)+2,        h=2^(2a-2).                       (9)
\end{verbatim}

These horizons are cofinal and satisfy (1). They have v2(t)=1 and binary
bit 20 equal to zero, since their only nonzero bits are 1 and
2a\textgreater=22. Put n0=delta*t, y0=Y\_n0(4t+1), a0=X\_n0(4t+1). The
congruence gives t\textbar y0, and the valuation gives v2(y0)=20. Hence

\begin{verbatim}
b0=y0/(delta*t)
\end{verbatim}

is an odd integer. Also a0=1 modulo 4 because n0 is even and the
abscissas at odd center are odd.

Use the inherited normalized Pell permutation lemma: if a0=1 modulo 4,
s-\textgreater Y\_s(a0) induces a permutation modulo every 2\^{}r. Here
is its unbounded proof. For H=2\^{}r, the valuation gives v2(Y\_H)=r.
The norm and X\_H=1 modulo 4 give X\_H=1 modulo 2\^{}(2r+2). Addition
then yields

\begin{verbatim}
Y_(s+2^r)(a0)-Y_s(a0)=2^r mod2^(r+1).
\end{verbatim}

The two lifts of each index class have opposite next output bit.
Induction from parity proves bijectivity and periodicity; adding
positive multiples of 2\^{}r gives cofinal indices in a selected class.

Choose b\textgreater22 with 2\^{}b\textgreater t. Since (t/2)b0 is odd,
the permutation supplies a positive odd indexs, as large as desired,
satisfying

\begin{verbatim}
(t/2)*b0*Y_s(a0)=1 mod2^(b-20).
\end{verbatim}

Set n=n0*s and

\begin{verbatim}
F=y0Y_s(a0)=Y_n(4t+1),
X=X_n(4t+1), Z=tF, u=(F/(delta*t)-1)/2.
\end{verbatim}

All values are integral, F is positive and arbitrarily large, and every
row of (2) holds. These choices give

\begin{verbatim}
F=2^20 mod2^b.                                      (10)
\end{verbatim}

Below bit b, F has only bit 20, which is absent from t; at or above bit
b, t has no nonzero bits. Therefore the binary supports of F and t are
disjoint. The identity v2(m!)=m-s2(m) gives

\begin{verbatim}
v2 binom(F+t,t)=s2(F)+s2(t)-s2(F+t)=0.
\end{verbatim}

Thus B is odd. This changes the incoming bit 23 construction to bit 20,
rather than silently reusing its old t=16 modulo 2\^{}26 horizon.

Choose an arbitrarily large positive odd w divisible by every odd prime
\textless= t, and put y=Fw. Every prime\textless=t divides y. If 1+iy
and 1+jy had a common prime p for 1\textless=i\textless j\textless=t+1,
then p is coprime to y and divides j-i, so p\textless=t; but then
p\textbar y, a contradiction. Hence all beta moduli are pairwise
coprime. This uses prime support, not the false assertion t!\textbar y
with fixed v2(F)=20. Choose w large enough to dominate all pending-state
values and then choose a positive beta code x arbitrarily large in its
CRT class. The odd-linear beta parameter is(w-1)/2. Define N=B\^{}x only
after x is fixed; N is odd.

The generalized factorial-free selector (8) assigns the independent CRT
representatives modulo N\^{}2. Choose the eight ordinary representatives
positive first. BecauseN\^{}2 is odd, chooseC even, positive and
cofinally large in its prescribed class, exceeding all other requested
indices and representatives. The selector congruenceB\textbar C+F+1
remains true. Only then setD\_radix, q and their exact powers. We
haveD\_radix\textgreater2C and N\textbar q, so the same digit bounds
supply all normalized quotients and remainders.

Choose the shared square frame after all these arithmetic outputs are
fixed. Its canonical positive principal pairs, private even-C anchor
comparisons, cap slack and orientation quotient are completed exactly by
the retained source construction. All their hypotheses have just been
checked, without a cofinal-square horizon. Evaluate every ordinary
circuit wire, substitute the new h in(1), omit the old root slot and
renumber. This constructs an integer zero from every finite refutation
and completes the ordinary equivalence.

\subsection{Exact verification and
portability}\label{exact-verification-and-portability}

\nolinkurl{semantic_certificate.json} pins the whole original210 equation,
every changed row, the exact substitution and the private root/row
deletion. \nolinkurl{frame_certificate.json} independently reconstructs
all 34 actual square atoms and covers all 45 actual power ports.
\nolinkurl{coordinate_map.json} gives all retained coordinates, rows and
exact renumbering. The independent checker imports neither producer nor
producer algebra library. It compares every row, verifies the inverse
affine polynomial identities on retained physical tuples, recollects
every quartic coefficient and parses all 298 plain/TeX residual displays
back into polynomials. Actual-data corruption controls must fail with
and without Python -O.

Verification is read-only unless an explicit receipt path is supplied.
Every verification rebuilds from a fresh relocated copy and
byte-compares all generated JSON. The optional parent replay
authenticates the original archive, the immediate principal210 stage,
its inherited source 173 closure, all source 210 stages and the full
original package verifier. New certificate paths are relative;
historical strings inside the unchanged original archive are provenance,
never paths used to find new inputs. No root search, remote write or
proof-assistant rebuild is involved. Finite checks certify literal
algebra and source binding, not the unbounded consistency theorem by
testing.
